% Source-derived chapter 06: Splittings of perverse filtrations
% Original source: hitchin-fibrations-abelian-surfaces-pw
\chapter{Splittings of perverse filtrations}
\label{hitchin-fibrations-abelian-surfaces-pw-chapter-06}
\source{hitchin-fibrations-abelian-surfaces-pw}

\begin{remark}[Source section]
\label{hitchin-fibrations-abelian-surfaces-pw-chapter-06-source}
\source{hitchin-fibrations-abelian-surfaces-pw}
\source{hitchin-fibrations-abelian-surfaces-pw}
This chapter reproduces the corresponding section of the retrieved
LaTeX source, including its definitions, statements, examples, and
proofs. It has not yet been translated into Lean.
\notready
\end{remark}

% The source section title was: Splittings of perverse filtrations
\subsection{Overview}
In this section, we study splittings of the perverse filtration associated with a proper surjective morphism $\pi: X\to Y$ with $X$ and $Y$ nonsingular. As an application, we strengthen Theorem \ref{taut_abelian} by requiring that the decomposition 
\begin{equation}\label{decomp6}
\source{hitchin-fibrations-abelian-surfaces-pw}
H^\ast(\CM_{\beta, A}, \BQ) = \bigoplus_{i,d} \widetilde{G}_iH^d(\CM_{\beta, A}, \BQ).
\end{equation}
is induced by a ``Lefschetz class" via the mechanisms introduced in \cite{D, dC, dCM4}; see Theorem \ref{strengthened2.1}. As discussed in Remark \ref{remark1}, this is crucial in the study of the specialization morphism (\ref{sp}) in Section 4. 

\subsection{Lefschetz classes}
We consider the perverse filtration associated with $\pi: X \to Y$,
\begin{equation}\label{e50}
\source{hitchin-fibrations-abelian-surfaces-pw}
P_0H^\ast(X, \BQ) \subset P_1H^\ast(X, \BQ) \subset \dots \subset P_{2r}H^\ast(X, \BQ) =  H^\ast(X, \BQ),
\end{equation}
where $r$ is the defect of semismallness (\ref{defect}). The action of a class $\eta \in H^2(X, \BQ)$ on the cohomology $H^*(X, \BQ)$ via cup product satisfies
\begin{equation}\label{e51}
\source{hitchin-fibrations-abelian-surfaces-pw}
\eta: P_kH^i(X, \BQ) \to P_{k+2}H^{i+2}(X, \BQ),
\end{equation}
which further induces an action on 
\[
\BH = \bigoplus_{i,j\geq 0} \mathrm{Gr}_i^P H^{i+j}(X, \BQ)= \bigoplus_i \left( P_iH^{i+j}(X, \BQ)/P_{i-1}H^{i+j}(X, \BQ) \right).
\]

We call $\eta \in H^2(X, \BQ)$ a \emph{Lefschetz class} with respect to the morphism $\pi: X\to Y$, if its induced action on $\BH$ satisfies the hard Lefschetz condition in the sense of \cite[Assumption 2.3.1]{dC}, \emph{i.e.}, the actions
\[
\eta^k: \mathrm{Gr}^P_{r-k}H^*(X, \BQ) \xrightarrow{\simeq}  \mathrm{Gr}^P_{r+k}H^*(X, \BQ), \quad \forall k\geq 0,
\]
are isomorphisms. As a typical example, the relative hard Lefschetz theorem \cite{BBD} with respect to $\pi: X \to Y$ implies that a relative ample class is a Lefschetz class. The following lemma gives examples of Lefschetz classes other than relative ample classes.  

\begin{lem}
\source{hitchin-fibrations-abelian-surfaces-pw}
\notready\label{Lemma3.1}
\source{hitchin-fibrations-abelian-surfaces-pw}
Let $A$ be an abelian variety of dimension $m$. Any class $\eta \in H^2(A, \BQ)$ satisfying
\[
\eta^m \neq 0 
\]
is a Lefschetz class with respect to $\pi: A \to \mathrm{pt}$.
\end{lem}

\begin{proof}
The action of $\eta$ on $H^*(A, \BQ)$ is a nilpotent operator. It suffices to show that $\eta$ can be completed into an $\mathfrak{sl}_2$-action. This follows from \cite[Proposition 3.3]{LL} and the Jacobson--Morosov lemma, since the operator $\eta$ is contained in the Loiijenga--Lunts semi-simple Lie algebra $\mathfrak{g}_{\mathrm{tot}}(A)$.
\end{proof}

Recall the Hitchin fibration $h: \CM_{\mathrm{Dol}} \to \Lambda$. The following proposition is a numerical criterion for Lefschetz classes with respect to $h$.

\begin{prop}
\source{hitchin-fibrations-abelian-surfaces-pw}
\notready\label{prop3.2_Hitchin}
\source{hitchin-fibrations-abelian-surfaces-pw}
Let $F$ be a closed fiber of $h$. A class $\eta \in H^2(\CM_{\mathrm{Dol}}, \BQ)$ is a Lefschetz class with respect to $h$ if and only if
\begin{equation}\label{eqn50}
\source{hitchin-fibrations-abelian-surfaces-pw}
\eta^{\mathrm{dim(F)}}|_{F} \neq 0.
\end{equation}
\end{prop}

\begin{proof}
We denote by $\hat{\CM}_{\mathrm{Dol}}$ the moduli space of stable $\mathrm{PGL}_n$-Higgs bundles with
\[
\hat{h}: \hat{\CM}_{\mathrm{Dol}} \to \hat{\Lambda}
\]
the corresponding Hitchin fibration. By \cite{HT1, Markman}, we have
\[
H^2( \hat{\CM}_{\mathrm{Dol}}, \BQ ) = \BQ \cdot \alpha
\]
where $\alpha$ is a relative ample class, and therefore a Lefschetz class, with respect to $\hat{h}$. By the discussion in \cite[Section 2.4]{dCHM1}, we have an isomorpism
\begin{equation}\label{112233}
\source{hitchin-fibrations-abelian-surfaces-pw}
H^*(\CM_{\mathrm{Dol}}, \BQ) = H^*( \hat{\CM}_{\mathrm{Dol}}, \BQ ) \otimes H^*(\mathrm{Pic}^0(C), \BQ)
\end{equation}
satisfying that
\[
P_k H^*(\CM_{\mathrm{Dol}}, \BQ) = \bigoplus_{i+j=k} P_iH^*( \hat{\CM}_{\mathrm{Dol}}, \BQ ) \otimes H^j(\mathrm{Pic}^0(C), \BQ).
\]
Under the identification 
\[
H^2(\CM_{\mathrm{Dol}}, \BQ) = H^2( \hat{\CM}_{\mathrm{Dol}}, \BQ ) \oplus H^2(\mathrm{Pic}^0(C), \BQ) = \BQ \cdot \alpha  \oplus H^2(\mathrm{Pic}^0(C)),
\]
induced by (\ref{112233}), we can express any class $\eta \in H^2(\CM_{\mathrm{Dol}}, \BQ)$ as
\[
\eta = \mu \alpha + \xi, \quad \mu \in \BQ,~ \xi \in H^2(\mathrm{Pic}^0(C), \BQ).  
\]
So by Lemma \ref{Lemma3.1}, the class $\eta$ is Lefschetz if and only if $\mu \neq 0$ and $\xi^g\neq 0$. This is equivalent to the condition (\ref{eqn50}).
\end{proof}

\subsection{The first Deligne splitting}
Deligne defined in \cite{D} three splittings of the perverse filtration associated with $\pi: X\to Y$.\footnote{Without the assumption that $X$ is nonsingular, the discusion in this section works identically for splitting the perverse filtration on the intersection cohomology $\mathrm{IH}^*(X, \BQ)$.} Each splitting relies on the choice of a Lefschetz class
\[
\eta \in H^2(X, \BQ).
\]
Deligne's constructions were later extended and generalized in \cite{dC, dCM4}. In this section, we only consider the splitting given by the first construction of \cite{D} which we call \emph{the first Deligne splitting}. 

%We follow closely the treatment of \cite{dCM4}.

The first Deligne splitting 
\[
H^*(X,\BQ) = \bigoplus_{i}G_iH^*(X, \BQ)
\]
of the perverse filtration (\ref{e50}) is governed by the sub-vector spaces
\[
Q^k \subset P_kH^*(X, \BQ),\quad  0\leq k\leq r
\]
satisfying that
\[
G_kH^*(X, \BQ) = \bigoplus_i \eta^i Q^{k-2i}.
\]
We describe $Q^k$ following \cite[Section 2.7]{dCM4}.

By (\ref{e51}), we have
\[
\eta^j P_k H^*(X, \BQ) \subset P_{k+2j}H^*(X, \BQ).
\]
Let $\Phi_0$ be the composition of the morphisms
\[
\Phi_0: P_kH^*(X, \BQ) \xrightarrow{\eta^{r-k+1}} P_{2r-k+2}H^*(X, \BQ) \to \mathrm{Gr}_{2r-k+2}^PH^*(X, \BQ),
\]
where the first map is cup product and the second map is the projection to the graded piece. We obtain the sub-vector space
\[
\mathrm{Ker}(\Phi_0) \subset P_kH^*(X, \BQ).
\]
The morphisms $\Phi_m$ are defined inductively for $m\geq 0$ as the following:
\[
\Phi_m: \mathrm{Ker}(\Phi_{m-1}) \xrightarrow{\eta^{r-k+m}} P_{2r-k+2m}H^*(X, \BQ) \to \mathrm{Gr}_{2r-k+2m}^PH^*(X, \BQ).
\]
Therefore for fixed $k$, we obtain a sequence of sub-vector spaces
\[
\cdots \subset \mathrm{Ker}(\Phi_1) \subset \mathrm{Ker}(\Phi_0) \subset P_kH^*(X, \BQ), 
\]
and $Q^k$ is obtained as
\[
Q^k = \mathrm{Ker}(\Phi_k) \subset P_kH^*(X, \BQ)
\]
by \cite[Proposition 2.7.1]{dCM4}.

The description of the first Deligne splitting yields immediately the following comparison lemma.

\begin{lem}
\source{hitchin-fibrations-abelian-surfaces-pw}
\notready\label{comparison1}
\source{hitchin-fibrations-abelian-surfaces-pw}
Let $X_i$ $(i=1,2)$ be nonsingular varieties with proper surjective morphisms $f_i: X_i \to Y_i$, and let $P_*H^*(X_i, \BQ)$ be the corresponding perverse filtrations. Assume there is a graded morphism 
\[
\phi: H^*(X_1, \BQ) \to H^*(X_2, \BQ)
\]
of $\BQ$-algebras satisfying 
\begin{enumerate}
    \item[(a)] $\phi(P_kH^i(X_1, \BQ)) \subset P_kH^i(X_2, \BQ)$;
    \item[(b)] $\phi(\eta_1) = \eta_2$ with $\eta_i \in H^2(X_i, \BQ)$ a Lefschetz class associated with $f_i$. 
\end{enumerate}
Then we have
\[
\phi(G_kH^i(X_1, \BQ)) \subset G_k H^i(X_2, \BQ), \quad \forall i,k,
\]
with $G_*H^*(X_i, \BQ)$ the first Deligne splitting associated with $\eta_i$.
\end{lem}

\begin{proof}
It follows from the fact that the sub-vector spaces
\[
\mathrm{Ker}(\Phi_m), ~ Q^k \subset P_kH^*(X_i, \BQ)
\]
are preserved under the morphism $\phi$.
\end{proof}

\begin{rmk}
\source{hitchin-fibrations-abelian-surfaces-pw}
\notready
In the case of the Hitchin fibration 
\[
h: \CM_{\mathrm{Dol}} \to  \Lambda,
\]
we see in the proof of Proposition \ref{prop3.2_Hitchin} that the first Deligne splitting associated with any Lefschetz class has the form 
\begin{equation}\label{Hitchin_splitting}
\source{hitchin-fibrations-abelian-surfaces-pw}
G_k H^*(\CM_{\mathrm{Dol}}, \BQ) = \bigoplus_{i+j=k} G_iH^*(\hat{\CM}_{\mathrm{Dol}}, \BQ) \otimes H^j(\mathrm{Pic}^0(C), \BQ), 
\end{equation}
where $G_iH^*(\hat{\CM}_{\mathrm{Dol}}, \BQ)$ is the first Deligne splitting associated with the relative ample class 
\[
\alpha \in H^2(\hat{\CM}_{\mathrm{Dol}}, \BQ).
\]
In particular, the splitting (\ref{Hitchin_splitting}) does not depend on the choice of a Lefschetz class.
\end{rmk}

\subsection{Semismall maps and Hilbert schemes}\label{Section3.4}
\source{hitchin-fibrations-abelian-surfaces-pw}
In this section, we study the morphism $\pi: X \to Y$ which can be factorized as
\[
X \xrightarrow{f} Z \xrightarrow{g} Y
\]
where $f: X \to Z$ is semismall. We further assume that there are sub-varieties $Z_i \subset Z$ such that the decomposition theorem for $f$ has the form
\begin{equation}\label{decomp9}
\source{hitchin-fibrations-abelian-surfaces-pw}
\mathrm{Rf_\ast} \BQ_X = \bigoplus_i \mathrm{IC}_{Z_i}[-\mathrm{dim}(X)].
\end{equation}
In particular, one of $Z_i$ is the total variety $Z$. The restriction of $g$ to every $Z_i$ gives the morphism
\[
g_i: Z_i \to Y_i \subset Y.
\]
We deduce from (\ref{decomp9}) the canonical decomposition of the cohomology of $X$:
\begin{equation}\label{decomp8}
\source{hitchin-fibrations-abelian-surfaces-pw}
H^d(X, \BQ) = \bigoplus_i \mathrm{IH}^{d-2n_i}(Z_i, \BQ), \quad n_i = \mathrm{codim}_{Z_i}(X).
\end{equation}

\begin{prop}
\source{hitchin-fibrations-abelian-surfaces-pw}
\notready\label{proposition3.5}
\source{hitchin-fibrations-abelian-surfaces-pw}
For $\alpha \in H^2(Z, \BQ)$, we assume that the restriction
\[
\alpha_i = \alpha|_{Z_i} \in H^2(Z_i, \BQ)
\]
is a Lefschetz class with respect to $g_i$ which induces the first Delgine splitting 
\[
\mathrm{IH}^*(Z_i, \BQ) = \bigoplus_k G_k \mathrm{IH}^*(Z_i, \BQ).
\]
Then $\eta = f^*\alpha \in H^2(X, \BQ)$ is a Lefschetz class, whose induced first Deligne splitting is
\begin{equation}\label{decomp7}
\source{hitchin-fibrations-abelian-surfaces-pw}
G_kH^d(X, \BQ) = \bigoplus_i G_{k-n_i}\mathrm{IH}^{d-2n_i}(Z_i, \BQ).
\end{equation}
under the identification (\ref{decomp8}).
\end{prop}

\begin{proof}
Since $f$ is semismall, it is a generically finite morphism. By the projection formula, we have
\[
f_\ast \eta = f_*f^* \alpha = \mathrm{deg}(f)\cdot  \alpha,
\]
which implies that the induced morphism
\[
f_\ast \eta: \bigoplus_i \mathrm{IC}_{Z_i} \to \bigoplus_i\mathrm{IC}_{Z_i}[2]
\]
is the direct sum 
\[
f_\ast \eta = f_*f^* \alpha = \mathrm{deg}(f) \cdot \bigoplus_i \alpha_i,\quad \alpha_i: \mathrm{IC}_{Z_i} \to \mathrm{IC}_{Z_i}[2].
\]
In particular, for a class written in the form
\[
(\gamma_i)_i \in H^*(X, \BQ), \quad \gamma_i \in \mathrm{IH}^{*}(Z_i, \BQ),
\]
via (\ref{decomp8}), we have $\eta \cup (\gamma_i)_i = \mathrm{deg}(f)\cdot(\alpha_i \cup \gamma_i)_i$.

We apply the decomposition theorem to the composition
\[
X \to Z \to Y
\]
and obtain that the perverse filtration $P_*H^*(X, \BQ)$ can be expressed in terms of the perverse filtrations $P_*\mathrm{IH}^*(Z_i, \BQ)$ under (\ref{decomp8}), \emph{i.e.},
\[
P_kH^d(X, \BQ) = \bigoplus_i P_{k-n_i}\mathrm{IH}^{d-2n_i}(Z_i, \BQ).
\]
Hence we deduce that $\eta$ is a Lefschetz class from the fact that every $\alpha_i$ is a Lefschetz class, and the decomposition (\ref{decomp7}) follows.
\end{proof}

We show in the following that the splitting (\ref{canonical_split}) for a Hilbert scheme is the first Deligne splitting induced by a Lefschetz class. 

The morphism $p_n: A^{[n]} \to E'^{(n)}$ introduced in Section \ref{Section1.5} can be factorized as
\begin{equation}\label{factorization_Hilb}
\source{hitchin-fibrations-abelian-surfaces-pw}
A^{[n]} \xrightarrow{f} A^{(n)} \xrightarrow{g} E'^{(n)}
\end{equation}
where the Hilbert--Chow morphism $f$ is semismall \cite{Go2}. There are canonical morphisms
\[
\kappa_\nu: A^{(\nu)} \to A^{(n)}
\]
together with a canonical stratification of $A^{(n)}$ indexed by the partitions $\nu$ of $n$,
\[
A^{(n)} = \bigcup_\nu A_\nu, \quad \overline{A_\nu} = \mathrm{Im}(\kappa_\nu) \subset A^{(n)}.
\]
By \cite[Theorem 3]{Go2}, the decomposition theorem associated with $f$ is
\[
\mathrm{Rf}_* \BQ_{A^{[n]}} = \bigoplus_{\nu} \mathrm{IC}_{\overline{A_\nu}} [-2n],
\]
and the restriction of $g: A^{(n)} \to E'^{(n)}$ to each $\overline{A_\nu}$ is described by the commutative diagrams
\begin{equation}\label{diagram9}
\source{hitchin-fibrations-abelian-surfaces-pw}
\begin{tikzcd}
A^{(\nu)} \arrow{r}{\kappa_\nu}  \arrow[swap]{d} &  \overline{A_\nu}\arrow[r, hookrightarrow]  \arrow[swap]{d} &  A^{(n)}  \arrow{d} \\
{E'}^{(\nu)} \arrow{r} & \overline{{E'_\nu}} \arrow[r, hookrightarrow] & E'^{(n)}.
\end{tikzcd}
\end{equation}

We consider the relative ample class
\[
\alpha = \mathrm{pt} \otimes 1_{E'} \in H^2(E, \BQ)\otimes H^0(E', \BQ) \subset H^2(A, \BQ)
\]
of $p:A \to E'$ which induces a Lefschetz class 
\[
\alpha^{(n)} \in H^2(A^{(n)}, \BQ)
\]
with respect to $A^{(n)} \to E'^{(n)}$. It further induces a Lefschetz class
\[
\alpha^{(\nu)} \in H^2(A^{(\nu)}, \BQ)
\]
with respect to $A^{(\nu)} \to E'^{(\nu)}$ for every partition $\nu$ of $n$. By the diagram (\ref{diagram9}), the pullback of $\alpha^{(n)}$ to every $A^{(\nu)}$ via $\kappa_\nu$ coincides with $\alpha^{(\nu)}$. Hence we obtain the following corollary by applying Proposition \ref{proposition3.5} to the factorization (\ref{factorization_Hilb}).

\begin{cor}
\source{hitchin-fibrations-abelian-surfaces-pw}
\notready\label{cor3.6}
\source{hitchin-fibrations-abelian-surfaces-pw}
The decomposition (\ref{canonical_split}) is the first Deligne splitting induced by the Lefschetz class 
\[
\eta_{A^{[n]}} = f^*\alpha^{(n)} \in H^2(A^{[n]}, \BQ)
\]
with respect to $p_n: A^{[n]} \to E'^{(n)}$.
\end{cor}

\subsection{A strengthened version of Theorem \ref{taut_abelian}}\label{Section3.5}
\source{hitchin-fibrations-abelian-surfaces-pw}
We study the decomposition (\ref{decomp6}) constructed for Theorem \ref{taut_abelian}. In the special case (\ref{special}), we have
\[
\CM_{\beta, A} \cong A^{[n]} \times \hat{A}
\]
with the morphism $\pi_\beta: \CM_{\beta,A} \to B$ given by (\ref{HC2}). Therefore we deduce from Corollary \ref{cor3.6} that the decomposition (\ref{splitting2}) is the first Deligne splitting associated with the Lefschetz class
\[
\eta_{A,\beta} = \eta_{A^{[n]}} \otimes 1_{\hat{A}} + 1_{A^{[n]}}\otimes \left(1_E \otimes \mathrm{pt} \right) \in H^2(A^{[n]} \times \hat{A}, \BQ) =  H^2(\CM_{\beta, A}, \BQ).
\]

By Theorem \ref{thm2.6}, any pair $(A', \beta')$ with $\beta'^2=2n$ is perversely isomorphic to the special pair $(A, \beta)$ given by (\ref{special}). So there is a graded isomorphism
\[
\phi: H^*(\CM_{\beta, A}, \BQ) \xrightarrow{\simeq} H^*(\CM_{\beta', A'}, \BQ) 
\]
of $\BQ$-algebras preserving the corresponding perverse filtrations, and we obtain a decomposition (\ref{decomp6}) as
\begin{equation}\label{996}
\source{hitchin-fibrations-abelian-surfaces-pw}
\widetilde{G}_kH^d(\CM_{\beta',A'}, \BQ) = \phi(\widetilde{G}_kH^d(\CM_{\beta',A'}, \BQ));
\end{equation}
see the proof of Proposition \ref{prop2.5}. Hence 
\[
\eta_{A',\beta'} = \phi(\eta_{A,\beta}) \in H^2(\CM_{\beta',A'}, \BQ)
\]
is a Lefschetz class with respect to $\pi_{\beta'}$, and Lemma \ref{comparison1} implies that the decomposition (\ref{996}) is the first Deligne splitting associated with $\eta_{A',\beta'}$. 

This gives the following strengthened version of Theorem \ref{taut_abelian}.

\begin{thm}
\source{hitchin-fibrations-abelian-surfaces-pw}
\notready\label{strengthened2.1}
\source{hitchin-fibrations-abelian-surfaces-pw}
Theorem \ref{taut_abelian} holds for a splitting $\widetilde{G}_\ast H^\ast(\CM_{\beta,A}, \BQ)$ given by the first Deligne splitting associated with a Lefschetz class with respect to $\pi_\beta: \CM_{\beta,A} \to B$.
\end{thm}
